\subsection{Tensor product of prehilbertian spaces}

Let \(E_{1}\) and \(E_{2}\) be two prehilbertian spaces and let \(F = E_{1} \otimes E_{2}\) be the tensor product of the vector spaces \(E_{1}\) and \(E_{2}\) . Let \(x_{1} \in E_{1}\) and \(x_{2} \in E_{2}\) ; since the mapping \((y_{1}, y_{2}) \mapsto \langle x_{1}|y_{1} \rangle \langle x_{2}|y_{2} \rangle\) from \(E_{1} \times E_{2}\) into K is bilinear, there exists a linear form \(\phi_{x_{1},x_{2}}\) on \(E_{1} \otimes E_{2}\) such that

\[
\phi_ {x _ {1}, x _ {2}} (y _ {1} \otimes y _ {2}) = \langle x _ {1} | y _ {1} \rangle \langle x _ {2} | y _ {2} \rangle\tag{1}
\]

for \(y_{1} \in \mathrm{E}_{1}\) and \(y_{2} \in \mathrm{E}_{2}\) . Let \(z \in \mathrm{F}\) . The mapping \((x_{1}, x_{2}) \mapsto \overline{\phi_{x_{1},x_{2}}(z)}\) from \(\mathrm{E}_1 \times \mathrm{E}_2\) into \(\mathbf{K}\) is bilinear; this can be seen by writing \(z\) in the form \(z = \sum_{i=1}^{n} y_{i,1} \otimes y_{i,2}\) with \(y_{i,1} \in \mathrm{E}_1\) and \(y_{i,2} \in \mathrm{E}_2\) for \(1 \leqslant i \leqslant n\) . Then there exists a linear form \(\psi_z\) on \(\mathrm{F} = \mathrm{E}_1 \otimes \mathrm{E}_2\) such that

\[
\psi_ {z} (x _ {1} \otimes x _ {2}) = \overline {{\phi_ {x _ {1} , x _ {2}} (z)}} \quad (x _ {1} \in \mathrm{E} _ {1}, x _ {2} \in \mathrm{E} _ {2}).\tag{2}
\]

We put \(\Phi(z, t) = \psi_z(t)\) for \(z, t\) in F. We see immediately that \(\Phi\) is a sesquilinear form on \(\mathbf{E}_1 \otimes \mathbf{E}_2\) characterized by

\[
\Phi (x _ {1} \otimes x _ {2}, y _ {1} \otimes y _ {2}) = \langle x _ {1} | y _ {1} \rangle \langle x _ {2} | y _ {2} \rangle\tag{3}
\]

(cf.~A, IX, § 1, No.~11).

PROPOSITION 1. --- The sesquilinear form \(\Phi\) on \(\mathrm{E}_1 \otimes \mathrm{E}_2\) is hermitian and positive, hence assigns the structure of a prehilbertian space to \(\mathrm{E}_1 \otimes \mathrm{E}_2\) . This space is Hausdorff if \(\mathrm{E}_1\) and \(\mathrm{E}_2\) are Hausdorff.

The formula \(\Phi(z, t) = \overline{\Phi(t, z)}\) follows from (3) when \(z = x_1 \otimes x_2\) and \(t = y_1 \otimes y_2\) . The general case is obtained by linearity, hence \(\Phi\) is hermitian.

Suppose \(\mathbf{E}_1\) and \(\mathbf{E}_2\) are Hausdorff; we shall prove that the hermitian form \(\Phi\) is positive and separating. Let \(z = \sum_{i=1}^{n} x_i \otimes y_i\) be a non-zero element of \(\mathrm{F} = \mathrm{E}_1 \otimes \mathrm{E}_2\) . Let \((e_1, ..., e_n)\) be an orthonormal basis of the subspace of \(\mathbf{E}_1\) generated by \(x_1, ..., x_n\) (V, p.~23, cor. 1). There exist elements \(f_1, ..., f_m\) in \(\mathbf{E}_2\) , not all null, such that \(z = \sum_{i=1}^{n} e_i \otimes f_i\) , hence

\[
\begin{array}{l} \Phi (z, z) = \sum_ {i, j = 1} ^ {m} \Phi (e _ {i} \otimes f _ {i}, e _ {j} \otimes f _ {j}) \\ = \sum_ {i, j} \langle e _ {i} | e _ {j} \rangle \langle f _ {i} | f _ {j} \rangle = \sum_ {i = 1} ^ {m} \| f _ {i} \| ^ {2} > 0. \end{array}
\]

For the general case, we shall now prove that \(\Phi\) is positive. Let \(\tilde{E}_{i}\) be the Hausdorff prehilbertian space associated with \(E_{i}\) and let \(\pi_{i}\) be the canonical mapping from \(E_{i}\) onto \(\tilde{E}_{i} (i = 1, 2)\) . Put \(\pi = \pi_{1} \otimes \pi_{2}\) . Let \(\tilde{\Phi}\) be the hermitian form on \(\tilde{E}_{1} \otimes \tilde{E}_{2}\) constructed in the same way as \(\Phi\) . We have

\[
\Phi (z, t) = \Phi (\pi (z), \pi (t)) \quad (z \in \mathrm{F}, t \in \mathrm{F}),
\]

and since \(\tilde{\Phi}\) is positive, so is \(\Phi\) .

Q.E.D.

The prehilbertian space defined in prop. 1 is called the tensor product of the prehilbertian spaces \(E_{1}\) and \(E_{2}\) and is written as \(E_{1} \otimes E_{2}\) . Henceforth we shall write \(\langle z|t\rangle\) for \(\Phi(z,t)\) , and therefore by definition

\[
\langle x _ {1} \otimes x _ {2} | y _ {1} \otimes y _ {2} \rangle = \langle x _ {1} | y _ {1} \rangle \langle x _ {2} | y _ {2} \rangle ;\tag{4}
\]

we shall also write \(\| z\| _2\) or \(\| z\|\) for \(\langle z|z\rangle^{1 / 2}\) . From (4), we get

\[
\left\| x _ {1} \otimes x _ {2} \right\| _ {2} = \left\| x _ {1} \right\|. \left\| x _ {2} \right\|,\tag{5}
\]

then the bilinear mapping \((x_{1}, x_{2}) \mapsto x_{1} \otimes x_{2}\) from \(\mathrm{E}_1 \times \mathrm{E}_2\) into \(\mathrm{E}_1 \otimes \mathrm{E}_2\) is continuous.

For i = 1, 2 let \(F_{i}\) be a vector subspace of \(E_{i}\) , endowed with the induced prehilbertian structure. Then \(F_{1} \otimes F_{2}\) can be identified with a vector subspace of \(E_{1} \otimes E_{2}\) (A, II, § 7, No.~7). Formula (4) shows that \(F_{1} \otimes F_{2}\) with the prehilbertian space structure induced by that of \(E_{1} \otimes_{2} E_{2}\) , is exactly \(F_{1} \otimes_{2} F_{2}\) . We shall henceforth identify \(F_{1} \otimes_{2} F_{2}\) with a prehilbertian subspace of \(E_{1} \otimes_{2} E_{2}\) .

PROPOSITION 2. --- For \(i = 1, 2\) , let \(\mathrm{E}_{i}\) and \(\mathrm{F}_{i}\) be two Hausdorff prehilbertian spaces and let \(u_{i} \in \mathcal{L}(\mathrm{E}_{i}; \mathrm{F}_{i})\) . The linear mapping \(u_{1} \otimes u_{2}\) from \(\mathrm{E}_{1} \otimes_{2} \mathrm{E}_{2}\) into \(\mathrm{F}_{1} \otimes_{2} \mathrm{F}_{2}\) is continuous and we have

\[
\left\| u _ {1} \otimes u _ {2} \right\| = \left\| u _ {1} \right\|. \left\| u _ {2} \right\|.
\]

Consider the positive hermitian form on \(\mathbf{E}_1\) given by

\[
f \left(x _ {1}, y _ {1}\right) = \| u _ {1} \| ^ {2} \left\langle x _ {1} \mid y _ {1} \right\rangle - \left\langle u _ {1} \left(x _ {1}\right) \mid u _ {1} \left(y _ {1}\right) \right\rangle .
\]

By prop. 1 (V, p.~25), there exists a positive hermitian form \(\Phi\) on \(E_1 \otimes E_2\) such that

\[
\begin{array}{r l} \Phi (x _ {1} \otimes x _ {2}, y _ {1} \otimes y _ {2}) & = f (x _ {1}, y _ {1}) \langle x _ {2} | y _ {2} \rangle = \\ & = \| u _ {1} \| ^ {2} \langle x _ {1} \otimes x _ {2} | y _ {1} \otimes y _ {2} \rangle - \langle (u _ {1} \otimes 1) (x _ {1} \otimes x _ {2}) | (u _ {1} \otimes 1) (y _ {1} \otimes y _ {2}) \rangle \end{array}
\]

for \(x_{1}, y_{1}\) in \(E_{1}\) and \(x_{2}, y_{2}\) in \(E_{2}\) . By linearity we have

\[
\Phi (z, t) = \| u _ {1} \| ^ {2} \langle z | t \rangle - \langle (u _ {1} \otimes 1) (z) | (u _ {1} \otimes 1) (t) \rangle
\]

for \(z, t\) in \(\mathrm{E}_1 \otimes \mathrm{E}_2\) . Since \(\Phi\) is positive, we get \(\Phi(z, z) \geqslant 0\) , that is, \(\| (u_1 \otimes 1).z \|_2\) \(\leqslant \| u_1\| \cdot \| z\| _2\) for \(z\in \mathrm{E}_1\otimes_2\mathrm{E}_2\) , or \(\| u_1\otimes 1\| \leqslant \| u_1\|\) . Similarly we prove the inequality \(\| 1\otimes u_2\| \leqslant \| u_2\|\) , and since \(u_{1}\otimes u_{2} = (u_{1}\otimes 1)\circ (1\otimes u_{2})\) , we get

\[
\left\| u _ {1} \otimes u _ {2} \right\| \leqslant \left\| u _ {1} \right\|. \left\| u _ {2} \right\|.
\]

On the other hand,

\[
\begin{array}{l} \| u _ {1} \| \cdot \| u _ {2} \| = \sup _ {\| x _ {1} \| \leqslant 1, \| x _ {2} \| \leqslant 1} \| u _ {1} (x _ {1}) \| \cdot \| u _ {2} (x _ {2}) \| \\ = \sup _ {\| x _ {1} \| \leqslant 1, \| x _ {2} \| \leqslant 1} \| (u _ {1} \otimes u _ {2}) (x _ {1} \otimes x _ {2}) \| _ {2} \leqslant \| u _ {1} \otimes u _ {2} \|. \end{array}
\]

This completes the proof of prop. 2.

Q.E.D.

Let \(E_1, \ldots, E_n\) be prehilbertian spaces ( \(n \geqslant 2\) ). We define the tensor product \(E_1 \otimes_2 \ldots \otimes_2 E_n\) (also denoted by \(\bigotimes_{i=1}^{n} E_i\) ) by induction, by

\[
\mathrm{E} _ {1} \otimes_ {2} \dots \otimes_ {2} \mathrm{E} _ {n} = (\mathrm{E} _ {1} \otimes_ {2} \dots \otimes_ {2} \mathrm{E} _ {n - 1}) \otimes_ {2} \mathrm{E} _ {n}.
\]

Hence, by the definition of scalar product, we have

\[
\left\langle x _ {1} \otimes \dots \otimes x _ {n} \mid y _ {1} \otimes \dots \otimes y _ {n} \right\rangle = \prod_ {i = 1} ^ {n} \left\langle x _ {i} \mid y _ {i} \right\rangle ,\tag{6}
\]

and in particular \(^{1}\)

\[
\left\| x _ {1} \otimes \dots \otimes x _ {n} \right\| _ {2} = \left\| x _ {1} \right\| \dots \left\| x _ {n} \right\|,\tag{7}
\]

for \(x_{i},y_{i}\) in \(\mathrm{E}_i\) ( \(1\leqslant i\leqslant n\) ). If the \(\mathrm{E}_i\) are Hausdorff, then so is \(\mathrm{E}_1\otimes_2\dots \otimes_2\mathrm{E}_n\) .

Let \(F_{1}, \ldots, F_{n}\) be prehilbertian spaces and \(u_{i} \in \mathcal{L}(E_{i}; F_{i})\) for \(1 \leqslant i \leqslant n\) . By induction on n, prop. 2 implies that \(u_{1} \otimes \ldots \otimes u_{n}\) is a continuous linear mapping from \(E_{1} \otimes_{2} \ldots \otimes_{2} E_{n}\) into \(F_{1} \otimes_{2} \ldots \otimes_{2} F_{n}\) and that

\[
\left\| u _ {1} \otimes \dots \otimes u _ {n} \right\| = \left\| u _ {1} \right\| \dots \left\| u _ {n} \right\|.\tag{8}
\]

Let \(\sigma \in \mathfrak{S}_n\) be a permutation of the set \(\{1, 2, \dots, n\}\) . Because of (6), the linear mapping \(p_{\sigma}\) from \(\mathrm{E}_1 \otimes_2 \dots \otimes_2 \mathrm{E}_n\) onto \(\mathrm{E}_{\sigma^{-1}(1)} \otimes_2 \dots \otimes_2 \mathrm{E}_{\sigma^{-1}(n)}\) characterized by

\[
p _ {\sigma} \left(x _ {1} \otimes \dots \otimes x _ {n}\right) = x _ {\sigma^ {- 1} (1)} \otimes \dots \otimes x _ {\sigma^ {- 1} (n)}\tag{9}
\]

is a prehilbertian space isomorphism (« commutativity of tensor product »).

Similarly, consider a partition of \(\{1, 2, \dots , n\}\) into \(m\) consecutive intervals \(\mathrm{I}_{1}, \dots , \mathrm{I}_{m}\) with \(\mathrm{I}_{k} = [a_{k}, a_{k + 1} - 1]\) for \(1 \leqslant k \leqslant m\). Put

\[
\mathrm{F} _ {k} = \bigotimes_ {i = a _ {k}} ^ {a _ {k + 1} - 1} \mathrm{E} _ {i} (1 \leqslant k \leqslant m).
\]

\(^{1}\) Here again we put \(\|z\|_{2} = \langle z|z\rangle^{1/2}\) for z in \(E_{1} \otimes_{2} \ldots \otimes_{2} E_{n}\) .

The canonical isomorphism from \(F_{1} \otimes \ldots \otimes F_{m}\) onto \(E_{1} \otimes \ldots \otimes E_{n}\) which transforms \(\bigotimes_{k=1}^{m} \bigotimes_{i=a_{k}}^{a_{k+1}-1} x_{i}\) into \(x_{1} \otimes \ldots \otimes x_{n}\) (A, II, §3, No.~9) is a prehilbertian space isomorphism (« associativity of the tensor product »).

